\section{Appendix B: Counting Theorems and Proof Sketches}

\subsection{\texorpdfstring{$\card{X_m}=F_{m+2}$}{|X\_m| = F\_{m+2}}}

Let $a_m=\card{X_m}$. Classify by the last bit: if $w_m=0$, the prefix can be any element of $X_{m-1}$; if $w_m=1$, then $w_{m-1}=0$, and the prefix corresponds to $X_{m-2}$. Thus
\[
a_m=a_{m-1}+a_{m-2},\quad a_1=2,\ a_2=3,
\]
which yields $a_m=F_{m+2}$.

\subsection{Cycle/boundary decomposition}

The boundary condition is equivalent to $w_1=w_m=1$. From the prohibition law, $w_2=w_{m-1}=0$, and the remaining middle segment of length $m-4$ has freedom equivalent to counting admissible words of length $m-4$, hence
\[
\card{X_m^{\mathrm{bdry}}}=\card{X_{m-4}}=F_{m-2}.
\]
The cyclic admissible set is the complement, giving
\[
\card{X_m^{\mathrm{cyc}}}=F_{m+2}-F_{m-2}=F_{m-1}+F_{m+1}.
\]
Simultaneously, $F_{m-1}+F_{m+1}$ is the Lucas number $L_m$ and satisfies $L_m=\Tr(A^m)$, where $A=\begin{pmatrix}1&1\\1&0\end{pmatrix}$.
